\chapter{Design of the Harness}
\label{ch:design}

This chapter presents the design: the target architecture, the ten principles the implementation follows, and the
separation between the harness and the evaluation code. Chapter~\ref{ch:impl} then describes how each part is built.

\section{Architecture}

Figure~\ref{fig:architecture} shows the target pipeline. A question first goes to a router, a cheap model or
classifier that picks the task family and the matching skill and extracts the parameters (node identifiers). The
agent loop then lets the cheap model solve the question with the tools of that family. A verifier written in plain
Python checks the final answer; if it passes, the answer is returned, and if it fails, the run is escalated to a
stronger model, which reruns the loop.

The layers, from the bottom up, are: the execution layer (networkx and, for large graphs, faster graph libraries); the
tool layer (validated graph tools, the result-handle store, and the MCP server); the model layer (LiteLLM
\cite{litellm}, which gives one interface to local and hosted models and reports token usage and cost); the harness
core (router, loop, verifiers, skills, escalation); and evaluation and logging on top.

On the date of this draft the agent loop, the tool layer, the verifiers, the model layer and the logging and
evaluation code are built. The router, the skills, escalation and the benchmark adapters are not; in the figure they
are drawn with dashed outlines.

\begin{figure}[ht]
  \centering
  \begin{tikzpicture}[
      node distance=6mm and 6mm,
      box/.style={draw, rounded corners, align=center, minimum height=10mm, text width=24mm, font=\small},
      planned/.style={box, dashed, text=black!70},
      arrow/.style={-{Stealth[length=2mm]}, thick}
    ]
    \node[box] (q) {question\\(parameters only)};
    \node[planned, right=of q] (router) {Router\\(task family,\\skill)};
    \node[box, right=of router] (loop) {Agent loop\\(cheap LLM $\leftrightarrow$\\graph tools)};
    \node[box, right=of loop] (verify) {Verifier\\(plain Python)};
    \node[box, below=12mm of verify] (answer) {answer\\(verified or\\unverified)};
    \node[planned, below=12mm of loop] (esc) {Escalation\\(stronger model\\reruns the loop)};
    \node[box, above=10mm of loop] (tools) {graph tools, handles\\(graph in memory)};
    \node[box, above=10mm of router] (llm) {LiteLLM\\(any model)};

    \draw[arrow] (q) -- (router);
    \draw[arrow] (router) -- (loop);
    \draw[arrow] (loop) -- (verify);
    \draw[arrow] (verify) -- node[right, font=\footnotesize] {pass} (answer);
    \draw[arrow] (verify) -- node[above, sloped, font=\footnotesize] {fail} (esc);
    \draw[arrow] (esc) -- (loop);
    \draw[arrow, <->] (loop) -- (tools);
    \draw[arrow, <->] (llm) -- (loop.north west);
  \end{tikzpicture}
  \caption[Target architecture of the harness]{Target architecture. Solid boxes are built; dashed boxes (router,
  escalation) are planned for M4 and M5. Today the evaluation runner picks the task directly from the dataset, which
  plays the router's role, and a failed verification is sent back to the same model as an error message.}
  \label{fig:architecture}
\end{figure}

A second view, the path of one question through the code as built, is given in Section~\ref{sec:one-question}.

\section{Design principles}

The implementation follows ten principles, fixed at the start of the project. Each is stated below with its
rationale and the way it is realised; deviations are noted.

\paragraph{1. The graph never goes into the prompt.}
The graph lives in Python memory; the model sees only tool results. This removes the dependence of the prompt size on
the graph size and moves the error-prone work (search, counting, copying) into code. For the dataset this means that
questions are rebuilt from a template with only the parameters, because the stored dataset questions contain the full
edge list.

\paragraph{2. Tools never raise.}
Every tool validates its inputs (the node exists, the graph has the right kind, the arguments have the right types)
and returns \texttt{\{"error": ...\}} with a clear, actionable message instead of raising. A crash would end the run;
an error message lets the model correct itself. In an early version, a tool that returned an object instead of a
number crashed the whole run; with this principle in place such bugs become an error message the model can read.

\paragraph{3. Result handles for compaction.}
Large outputs are stored on the harness side, and the model receives a short summary plus an identifier, for example
\texttt{\{"handle": "result\_1", "length": 9932, "first\_5": [...]\}}. Tools and \code{submit_answer} accept handles, so a
9,932-edge spanning forest can be submitted without the model ever writing it out (Section~\ref{sec:handles}). The
principle was later generalised: nothing the model reads may grow with the graph, including error messages and
repeated values.

\paragraph{4. Structured final answers.}
A run ends when the model calls \code{submit_answer}, whose schema depends on the task (a number, true/false, a node
list, an edge list, or a yes/no with evidence). A shape error goes back to the model. The schema matters more than the
instructions: in one experiment, the model was told it could submit a handle, but the schema said ``array'', and it
did the best the schema allowed (Section~\ref{sec:handles}). A second ending tool, \code{cannot_answer}, gives the
model an honest exit instead of a guess.

\paragraph{5. Verification is code, not the model, and checks evidence.}
Each task family has a programmatic verifier. A verifier may check only what the answer itself shows: that a path is
real and connects the right nodes, that a cycle is a cycle, that an edge set is a spanning forest. It must not
recompute the answer and compare it with the submission, because that would be an oracle: accuracy would approach
100\% and any comparison with another harness would be meaningless. Answers that cannot be checked without
re-solving (counts, degrees, edge existence, and every ``no'') are run without a verifier and reported as unverified.
Table~\ref{tab:checkable} summarises the resulting split.

One qualification belongs here. The shortest-path verifier and the waypoint-route verifier check the path against the
graph and then compare its length with the shortest-path distance computed by breadth-first search. The path itself is
evidence, but the optimality check does compute the optimal length. A certificate-based check, in which distance
labels that satisfy every edge prove a length optimal in $O(|E|)$ checks without a search, is noted as a possible
upgrade (\code{docs/gds_agent_findings.md}).
\todo[inline]{Decide whether the BFS length check in verify\_shortest\_path / verify\_path\_via is acceptable under principle 5, or replace it with a distance-label certificate; state the decision here.}

\begin{table}[ht]
  \centering
  \small
  \caption{What can be checked from evidence alone, per task.}
  \label{tab:checkable}
  % source: docs/learning_log.md (2026-10-06), harness/tasks.py, harness/verifiers.py
  \begin{tabularx}{\textwidth}{>{\raggedright\arraybackslash}p{5cm}>{\raggedright\arraybackslash}X}
    \toprule
    Task & Check \\
    \midrule
    \code{shortest_path}, \code{mst}, \code{shortest_path_via} & Full: the answer is the evidence (real edges, right endpoints; no cycle and $n-c$ edges for a spanning forest) \\
    \code{connectivity} ``yes'', \code{cycle_check} ``yes'' & Partial: the model must send the path or cycle, which is checked; a ``no'' is accepted unseen \\
    \code{connected_nodes} & Partial: listed nodes must be real neighbors, each listed once; a missing neighbor cannot be seen \\
    counts, \code{node_degree}, \code{edge_existence}, any ``no'', combine tasks & None: unverified by design \\
    \bottomrule
  \end{tabularx}
\end{table}

\paragraph{6. Show only relevant tools.}
The router picks a task family and the agent sees only that family's tools. Every tool on offer adds to every prompt
and is one more option to choose wrongly. This principle is not yet realised: the router is planned for M4, and today
the model sees all default tools. The cost is visible in the data: adding three tools in M3 grew the tool list the
model reads on every call from about 2,000 to about 2,800 characters.

\paragraph{7. Budgets everywhere.}
Runs have a maximum number of steps (10 by default), a loop detector that stops a run when the model sends the same
reply three times in a row, and, since the context-window problem described in Section~\ref{sec:case-4k}, a stop when
the conversation no longer fits. A maximum cost is part of the principle but is not needed for local models, whose
dollar cost is zero.
\todo[inline]{A per-run dollar budget is not implemented yet; add it before API models are used (M5/M6).}

\paragraph{8. Append-only message history.}
Messages are only appended, which keeps provider prompt caches effective (the prefix of every call is the previous
call). Two deliberate exceptions exist: the model's hidden thinking is removed from its reply before the reply joins
the history (it was otherwise re-sent and re-read on every step), and compaction shortens older tool results near the
context limit, which restarts the cache once.

\paragraph{9. Model-agnostic.}
All model calls go through LiteLLM; no provider is hard-coded. Switching model is changing one string, and the same
loop runs local Ollama models and hosted API models.

\paragraph{10. Log everything.}
Every run writes a JSONL trace: messages, tool calls, results as the model saw them, tokens, latency and, for hosted
models, cost. Every number in Chapter~\ref{ch:results} is reproducible from these files.

\section{Separation of harness and evaluation}

The harness (\code{harness/}) and the evaluation code (\code{eval/}) are separated in one direction: evaluation may
import the harness, never the reverse. The task registry in the harness holds only what the agent may use: the
question template, the answer type and the verifier. Reference answers, dataset parsing and grading live in
\code{eval/tasks.py}, and the grader is written independently of the verifiers. The sandbox that runs model-written
code mounts only the graph file and the harness package, never \code{eval/} or \code{results/}, which hold reference
answers. This is a matter of experimental validity as much as of safety: no ground truth can leak into a run.

A second split concerns the data: graphs 0--49 of each size are the development split, used to build and tune tools,
prompts and routing, and graphs 50--99 are the test split, kept untouched for final results. The split is by graph,
so no test graph is ever seen during development.
